% !TeX root = ../proyecto.tex
%--------------------------------------
\section{IU-06 Panel de Gráficas y Estadísticas}

\subsection{Objetivo}
	Proporcionar al analista una visión cuantitativa del corpus clasificado mediante cuatro
	visualizaciones interactivas que resumen la distribución temporal, temática, geográfica
	y por fuente del conjunto de noticias procesadas.

\subsection{Diseño}
	Este panel forma parte de la sección inferior del \IUref{IU-01}{Dashboard Principal},
	presentándose como una cuadrícula de cuatro tarjetas de gráfica que se renderizan con datos
	obtenidos en tiempo real de la API. Cada tarjeta tiene un título descriptivo, la gráfica
	correspondiente y un indicador de la fuente de datos activa (\texttt{postgresql} o
	\texttt{csv}). Los colores de todas las gráficas siguen el esquema cromático del sistema
	(escala de rojos y naranjas para relevancia Alta, grises para irrelevante).

\IUfig[.9]{panel_estadisticas}{IU-06}{Panel de Gráficas y Estadísticas del Corpus.}

\subsection{Gráficas}
	\begin{itemize}
		\item \textbf{Distribución Temporal} (\texttt{/api/charts/temporal}): Gráfica de línea o
		      barras agrupadas que muestra el número de noticias recolectadas por mes, con una
		      línea secundaria para las noticias con NNA identificados. Permite observar la
		      evolución del corpus a lo largo del tiempo.

		\item \textbf{Distribución de Relevancia} (\texttt{/api/charts/relevancia}): Gráfica de
		      dona (\textit{donut chart}) que presenta las proporciones de noticias clasificadas
		      como Alta, Media, Baja y No relevante sobre el total del corpus activo.

		\item \textbf{Distribución Geográfica} (\texttt{/api/charts/estados}): Gráfica de barras
		      horizontales con los 15 estados de México con mayor frecuencia de menciones en el
		      corpus, calculada mediante detección de patrones de nombres de estados y ciudades
		      clave en el contenido de las noticias.

		\item \textbf{Top de Fuentes} (\texttt{/api/charts/fuentes}): Gráfica de barras verticales
		      con las 10 fuentes periodísticas que más noticias aportaron al corpus en el lote
		      activo, facilitando la evaluación de la cobertura del catálogo.
	\end{itemize}

\subsection{Salidas}
	\begin{itemize}
		\item Las cuatro gráficas renderizadas dinámicamente con los datos del lote seleccionado.
		\item Indicador de fuente de datos activa en cada tarjeta.
		\item Estado de carga (spinner) mientras se obtienen los datos de la API.
	\end{itemize}

\subsection{Entradas}
	\begin{itemize}
		\item Selector de lote de análisis (\texttt{batch\_id}) heredado del \IUref{IU-01}{Dashboard}
		      que actualiza todas las gráficas de forma sincronizada.
	\end{itemize}

\subsection{Comandos}
	\begin{itemize}
		\item \IUbutton{Actualizar Gráficas}: Solicita nuevamente los datos a los cuatro endpoints
		      de gráficas para refrescar las visualizaciones con los datos más recientes.
	\end{itemize}
